\hypertarget{high-level-url-and-http-handling-urllib-http}{%
\chapter{\texorpdfstring{High-Level URL and HTTP Handling
(\texttt{urllib},
\texttt{http})}{High-Level URL and HTTP Handling (urllib, http)}}\label{high-level-url-and-http-handling-urllib-http}}

While low-level sockets (Chapter 38) are for systems plumbing, most
application-level networking uses HTTP. Python provides a layered suite
of modules to handle URLs and the HTTP protocol state machine.

\hypertarget{urllib.parse-the-url-state-machine}{%
\subsection{\texorpdfstring{54.1 \texttt{urllib.parse}: The URL State
Machine}{54.1 urllib.parse: The URL State Machine}}\label{urllib.parse-the-url-state-machine}}

A URL is not just a string; it is a complex address with hierarchy and
parameters (RFC 3986). * \textbf{\passthrough{\lstinline!urlparse()!}}:
Breaks a string into a 6-item named tuple
(\passthrough{\lstinline!scheme!}, \passthrough{\lstinline!netloc!},
\passthrough{\lstinline!path!}, \passthrough{\lstinline!params!},
\passthrough{\lstinline!query!}, \passthrough{\lstinline!fragment!}). *
\textbf{Safety}: Modern CPython has hardened
\passthrough{\lstinline!urllib.parse!} to prevent ``domain name
splitting'' attacks where special Unicode characters (like
\passthrough{\lstinline!\\uff01!}) are used to trick servers into
misrouting requests.

\hypertarget{http.client-the-protocol-engine}{%
\subsection{\texorpdfstring{54.2 \texttt{http.client}: The Protocol
Engine}{54.2 http.client: The Protocol Engine}}\label{http.client-the-protocol-engine}}

\passthrough{\lstinline!http.client!} is the lowest level of HTTP
handling before raw sockets. * \textbf{Persistence}: It supports
HTTP/1.1 persistent connections
(\passthrough{\lstinline!Connection: keep-alive!}). *
\textbf{Streaming}: You can send and receive request bodies in chunks
using the \passthrough{\lstinline!chunked!} transfer encoding, which is
essential for uploading large files without consuming all system memory.

\hypertarget{urllib.request-the-opener-pipeline}{%
\subsection{\texorpdfstring{54.3 \texttt{urllib.request}: The Opener
Pipeline}{54.3 urllib.request: The Opener Pipeline}}\label{urllib.request-the-opener-pipeline}}

\passthrough{\lstinline!urllib.request!} provides a high-level API built
on an extensible ``Handler'' architecture. 1. \textbf{Handlers}: Objects
that handle specific schemes (\passthrough{\lstinline!HTTPHandler!},
\passthrough{\lstinline!FTPHandler!},
\passthrough{\lstinline!FileHandler!}). 2. \textbf{Opener}: The
\passthrough{\lstinline!OpenerDirector!} manages a list of handlers.
When you call \passthrough{\lstinline!urlopen()!}, it iterates through
handlers until one accepts the request. 3. \textbf{Hooks}: You can write
custom handlers to implement caching, authentication, or automatic retry
logic.

\begin{center}\rule{0.5\linewidth}{0.5pt}\end{center}
